% outlook/part_intro-DE.tex
\chapter*{Ausblick}
\phantomsection
\addcontentsline{toc}{chapter}{Ausblick}

Dieser Teil zeichnet den weiteren Weg von AIM nach Abschluss der
wissenschaftlichen Argumentation dieser Arbeit vor.

\section*{Was wissen wir bereits?}

Die vorangehenden Teile haben das AIM-Framework in seiner gegenwärtigen
Form entwickelt und analysiert.

Grundlagen, Geometrie, Zustände, Laufzeitdienste, Realisierungsbegriffe,
Modellierungsstrukturen, Optimierungsmethoden, Ingenieurbedingungen,
Anwendungen und weiter reichende Interpretationen sind nun entwickelt.

Zusammen bilden diese Elemente ein kohärentes Framework für das
Alignment-Based Information Modelling.

Die vorliegende Arbeit liefert damit eine vollständige Beschreibung von
AIM in seiner gegenwärtigen Form.

\section*{Was fehlt noch?}

Kein Framework ist jemals wirklich abgeschlossen.

Neue Technologien entstehen, Ingenieurarbeitsabläufe entwickeln sich
weiter, Informationssysteme werden zunehmend integriert, und aus
Automatisierung, digitalen Zwillingen, Infrastrukturmanagement sowie
künftigen Verkehrssystemen erwachsen neue Anforderungen.

Das noch fehlende Element ist daher die Zukunftsperspektive. Die Frage
lautet nicht länger, was AIM ist, sondern was AIM werden könnte.

\section*{Was folgt?}

Dieser Teil stellt AIM in einen weiteren Zusammenhang und untersucht
mögliche Richtungen seiner künftigen Entwicklung.

Die folgenden Kapitel behandeln:
\begin{itemize}
\item BIM-bezogene Perspektiven,
\item Alignment-zentrierte Informationsmodelle,
\item künftige Ingenieurarbeitsabläufe,
\item Erweiterungen in Richtung von PIM-Begriffen
\item und offene Forschungsfragen.
\end{itemize}

Das Ziel besteht nicht darin, das Framework endgültig abzuschließen,
sondern Möglichkeiten für Wachstum und Integration aufzuzeigen.

Besondere Aufmerksamkeit gilt der Möglichkeit, dass sich das Alignment
zu einer eigenständigen Informationseinheit entwickelt, vergleichbar mit
der Rolle, die heute gebäudezentrierte BIM-Begriffe einnehmen.

\section*{Bezug zur AIM-Roadmap}

Innerhalb der AIM-Roadmap entspricht dieser Teil dem Übergang vom
gegenwärtigen Framework zu künftigen Entwicklungen.

Das folgende Symbol behandelt die abschließende Orientierungsfrage der
Arbeit: Welche Kontinuitätslinie verbindet die etablierten Ergebnisse mit
der Forschungsfront?

\begin{figure}[ht]
\centering
% figures/overview/icon_outlook_established_frontier-DE.tex
\begin{tikzpicture}[
  node distance=1.3cm,
  box/.style={draw, rounded corners, minimum width=5.0cm, minimum height=0.9cm, align=center},
  secondary/.style={draw, rounded corners, dashed, minimum width=5.0cm, minimum height=0.9cm, align=center},
  arrow/.style={->, thick}
]
\node[box] (established) {Etabliertes Wissen};
\node[secondary, below=of established] (frontier) {Forschungsfront};
\draw[arrow] (established) -- (frontier);
\draw[thin] ($(established.center)+(-3.0,-1.3)$) -- ($(established.center)+(3.0,-1.3)$);
\end{tikzpicture}

\caption{Ausblickssymbol: Kontinuität vom etablierten Wissen zur Forschungsfront.}
\label{fig:icon-outlook-established-frontier}
\end{figure}

Die Entwicklung
\[
\text{gegenwärtiges AIM}
\rightarrow
\text{künftige Entwicklung}
\rightarrow
\text{BIM, PIM und darüber hinaus}
\]
kennzeichnet den Übergang von der Beschreibung des gegenwärtigen
Frameworks zur Erkundung seiner künftigen Entwicklung.

Der Ausblick dient daher nicht als Schluss, sondern als Einladung, das
Alignment-Based Information Modelling weiterzuentwickeln.

\begin{principle}[Prinzip des offenen Frameworks]
\label{princ:open-framework-principle}
Alignment-Based Information Modelling soll in AIM als sich entwickelndes
Framework verstanden werden, dessen künftige Entwicklung voraussichtlich
über die in dieser Arbeit vorgestellten Begriffe hinausreichen wird.
\end{principle}

\section*{Rückkehr zum Ausgangspunkt}

Die Thesis begann mit einem Problem des Eisenbahnübergangs. Berlinish stellte
eine verallgemeinerte funktionale Grammatik bereit: einlaufende Halbwelle,
optionaler Klothoidenkern und auslaufende Halbwelle einschließlich
Null-Längen- und asymmetrischer Fälle. axtranNew stellte die
Berechnungsabsicht bereit: den Anschluss lösen, seine Freiheitsgrade offenlegen,
zulässige Alternativen vergleichen und editierbare Alignment-Struktur
zurückgeben.

Die weitere Arbeit ergänzte, was diese beiden Beiträge allein nicht leisten
konnten. transitionDB und TransEd bewahren und untersuchen
wiederverwendbares Funktionswissen. SPOT bewahrt Identität und Zustand
konkreter Ingenieurobjekte. alignmentOS macht Alignment-Wissen operativ. Der
Knowledge Kernel stabilisiert Terminologie und Verantwortungsgrenzen. Research
prüft Aussagen und eröffnet Erweiterungen. Die Thesis macht aus der
Entwicklungslinie ein überprüfbares und übertragbares Argument.

Etabliert ist die in dieser Thesis entwickelte krümmungsbasierte Kompositions-
und Zustandsmaschinerie. Implementierungsbelegt sind die gegenwärtige Registry,
Rekonstruktion, Auswertung sowie ausgewählte Bearbeitungs- und
Optimierungspfade. Offen bleiben zuverlässige allgemeine Optimierung,
validierter getrennter oder gekoppelter Krümmungs--Überhöhungs-Entwurf,
umfassendere Netzberechnung und eine etwaige spätere Kernel-Behandlung der
ursprünglichen Berlinish-Terminologie. Die Zukunft beginnt an dieser
ausdrücklichen Grenze und nicht mit einer Behauptung von Vollständigkeit.

\begin{summary}
Die Diskussion erklärt, was AIM bedeutet.

Der Ausblick erkundet, was AIM werden könnte.

Dieser Teil benennt daher künftige Richtungen, weitere Anwendungen und
offene Fragen, die über den Umfang der gegenwärtigen Arbeit hinausgehen
und zugleich auf ihren Grundlagen aufbauen.
\end{summary}

Insgesamt fasst diese abschließende Perspektive die Arbeit als
durchgehendes Argument von der grundlegenden Motivation bis zur künftigen
Ingenieurentwicklung.
